\documentclass[10pt,a3paper]{article}
\usepackage[margin=10mm]{geometry}
\usepackage[T1]{fontenc}
\usepackage[utf8]{inputenc}
\usepackage[french]{babel}
\usepackage{lmodern}
\usepackage{microtype}
\makeatletter\def\input@path{{../../../Style/}}\makeatother
\NeedsTeXFormat{LaTeX2e}
\ProvidesPackage{TLPosterCarteMentale}[2026/08/03 Posters et cartes mentales SI]

\RequirePackage{tikz}
\RequirePackage[most]{tcolorbox}
\RequirePackage{xcolor}
\RequirePackage{amsmath,amssymb}
\RequirePackage{graphicx}
\usetikzlibrary{positioning,arrows.meta,calc,shapes.geometric,mindmap,backgrounds,fit}

\definecolor{TLPosterBlue}{RGB}{24,86,141}
\definecolor{TLPosterLight}{RGB}{234,243,250}
\definecolor{TLPosterAccent}{RGB}{232,126,34}
\definecolor{TLPosterGreen}{RGB}{52,152,102}
\definecolor{TLPosterRed}{RGB}{192,57,43}
\definecolor{TLPosterGray}{RGB}{80,90,100}

\newcommand{\TLPosterTitre}[2]{%
  \begin{tcolorbox}[enhanced,boxrule=0pt,arc=0pt,colback=TLPosterBlue,left=6mm,right=6mm,top=4mm,bottom=4mm]
    {\color{white}\bfseries\LARGE #1}\hfill{\color{white}\large #2}
  \end{tcolorbox}%
}

\newtcolorbox{TLPosterBloc}[2][]{
  enhanced,breakable,
  colback=TLPosterLight,
  colframe=TLPosterBlue,
  colbacktitle=TLPosterBlue,
  coltitle=white,
  fonttitle=\bfseries,
  title={#2},
  boxrule=0.8pt,
  arc=2mm,
  left=3mm,right=3mm,top=2mm,bottom=2mm,
  #1
}

\newtcolorbox{TLPosterFormule}[1][]{
  enhanced,
  colback=white,
  colframe=TLPosterAccent,
  boxrule=1pt,
  arc=2mm,
  fontupper=\bfseries,
  halign=center,
  #1
}

\newcommand{\TLPosterMethode}[1]{%
  \begin{tcolorbox}[enhanced,colback=TLPosterGreen!8,colframe=TLPosterGreen,title=Méthode,fonttitle=\bfseries]
  #1
  \end{tcolorbox}%
}

\newcommand{\TLPosterAttention}[1]{%
  \begin{tcolorbox}[enhanced,colback=TLPosterRed!6,colframe=TLPosterRed,title=Point de vigilance,fonttitle=\bfseries]
  #1
  \end{tcolorbox}%
}

\tikzset{
  TLmindroot/.style={concept,concept color=TLPosterBlue,text=white,font=\bfseries\large,minimum size=34mm},
  TLmindlevel1/.style={level 1 concept/.append style={font=\bfseries,minimum size=23mm,sibling angle=45}},
  TLmindlevel2/.style={level 2 concept/.append style={font=\small,minimum size=17mm}},
  TLarrow/.style={-{Stealth[length=3mm]},thick,draw=TLPosterBlue},
  TLbox/.style={rounded corners=2mm,draw=TLPosterBlue,fill=TLPosterLight,align=center,inner sep=3mm}
}

\newenvironment{TLCarteMentale}[1]{%
  \begin{tikzpicture}[mindmap,grow cyclic,concept color=TLPosterBlue,every node/.style={concept},TLmindlevel1,TLmindlevel2]
  \node[TLmindroot]{#1}
}{;
  \end{tikzpicture}
}

\newcommand{\TLBranche}[3][]{child[concept color=#2]{node[#1]{#3}}}

\endinput

\pagestyle{empty}
\renewcommand{\familydefault}{\sfdefault}
\begin{document}
\TLPosterCourseHeader{8}{Conversion électromécanique}{Ramener la charge au moteur et identifier le régime de fonctionnement}
\vspace{2mm}
\begin{minipage}[t]{0.60\linewidth}
\begin{TLPosterSavoirs}
\begin{itemize}
  \item utiliser $P=C\Omega$, rendements et rapports cinématiques pour ramener une charge à l'arbre moteur
  \item distinguer charge résistante et entraînante et étudier la réversibilité de toute la chaîne
  \item reconnaître les lois de charge : constante, $k\Omega$, $k/\Omega$ ou $k\Omega^2$
  \item déterminer point de fonctionnement, accélération, freinage et quadrants
  \item exploiter essais à vide, en charge et rotor bloqué selon le paramètre recherché
\end{itemize}
\end{TLPosterSavoirs}
\end{minipage}\hfill
\begin{minipage}[t]{0.37\linewidth}
\TLPosterKey{Quel calcul ?}{Régime établi $\rightarrow C_u=C_r$\\Accélération $\rightarrow J_{eq}\dot\Omega=C_u-C_r$\\Charge à ramener $\rightarrow$ rapport + bilan de puissance\\Mode moteur/génératrice $\rightarrow$ signe de $C\Omega$}
\end{minipage}
\vfill
\begin{minipage}[t]{0.49\linewidth}
\begin{TLPosterBloc}[Ramener la charge à l'arbre moteur]{TLPosterBlue}
Traiter chaque transmission comme un adaptateur. Pour $k=\Omega_s/\Omega_e$ :
\[
\Omega_s=k\Omega_e,\qquad C_s\Omega_s=\eta C_e\Omega_e.
\]
Les vitesses se ramènent avec les rapports ; les couples avec le \textbf{bilan de puissance et les rendements}. Pour plusieurs étages : $\eta_{glob}=\prod\eta_i$. Toujours vérifier si les frottements sont déjà inclus dans $C_r$.
\end{TLPosterBloc}
\begin{TLPosterBloc}[Point de fonctionnement et stabilité]{TLPosterPurple}
En régime établi :
\[
C_u(\Omega)=C_r(\Omega).
\]
Le point de fonctionnement est l'intersection machine--charge. Si la vitesse augmente légèrement, le couple résultant $C_u-C_r$ doit devenir négatif pour ramener le système vers l'équilibre : c'est le test de stabilité du point.
\end{TLPosterBloc}
\end{minipage}\hfill
\begin{minipage}[t]{0.49\linewidth}
\begin{TLPosterBloc}[Transitoire, quadrants et freinage]{TLPosterGreen}
Avec inertie équivalente :
\[
J_{eq}\dot\Omega=C_u-C_r\quad(\text{ou }C_u-C_r-f\Omega).
\]
Accélération : $C_u>C_r$ ; régime établi : égalité ; décélération : $C_u<C_r$. Le signe de $P_m=C\Omega$ donne le mode : $P_m>0$ moteur, $P_m<0$ génératrice/frein. Une récupération d'énergie exige une chaîne \textbf{entièrement réversible}.
\end{TLPosterBloc}
\begin{TLPosterBloc}[Quel essai utiliser ?]{TLPosterTeal}
\textbf{À vide :} $C_r=0$ $\Rightarrow$ estimation du couple de pertes et des pertes fer.\\
\textbf{En charge :} rendement, puissances et pertes proches du fonctionnement réel.\\
\textbf{Rotor bloqué :} $\Omega=0$ sous tension réduite $\Rightarrow$ identification de paramètres électromagnétiques.\\[1mm]
Comparer ensuite le point trouvé aux valeurs nominales de la plaque.
\end{TLPosterBloc}
\end{minipage}
\vfill
\begin{TLPosterMethode}
\begin{enumerate}
  \item Identifier la charge, le sens du mouvement et le sens possible du flux d'énergie.
  \item Ramener vitesse, couple et éventuellement inertie à l'arbre moteur ; inclure les rendements.
  \item Choisir le régime : établi $\Rightarrow C_u=C_r$ ; transitoire $\Rightarrow$ PFD sur l'arbre moteur.
  \item Déterminer le point de fonctionnement ou le couple nécessaire sur chaque phase du profil de vitesse.
  \item Utiliser le signe de $C\Omega$ pour identifier le quadrant et le besoin de freinage dissipatif ou récupératif.
  \item Calculer puissances et rendement, puis vérifier courant, couple, vitesse et fonctionnement nominal.
\end{enumerate}
\end{TLPosterMethode}
\vfill
\TLPosterRetenir{avant de choisir une machine, tout ramener à son arbre : vitesse, couple, inertie et puissance ; puis étudier chaque phase du cycle et chaque quadrant demandé.}
\end{document}
